\documentclass[sigconf,nonacm=false]{acmart}
\usepackage{booktabs}
\usepackage{graphicx}
\usepackage{amsmath}
\usepackage{xcolor}
\usepackage{tikz}
\usepackage[ruled,vlined]{algorithm2e}

\begin{document}
\title{On-Chain AI Agent Identification under Label Leakage: An Honest Empirical Study with Provenance-Only Ground Truth}
\author{NSF Project Team}
\affiliation{\institution{University}\country{USA}}

\begin{abstract}
Autonomous AI agents manage billions on Ethereum but are invisible to on-chain monitoring. We build a 23-feature behavioral classifier on 1,147 provenance-verified addresses and discover that our own labeling heuristic inflated AUC from 0.98 to an honest 0.80. Most surprisingly, the widely assumed ``agents are riskier'' narrative reverses under honest labels: human DeFi users hold more unlimited token approvals (66.1 mean) than curated MEV bots (49.1 mean), because heuristic labels systematically misclassified heavy DeFi users as agents. Temporal holdout reveals random CV overestimates AUC by 13--16 points; the deployment-realistic AUC is 0.67 with precision 0.81 but recall 0.31, meaning the classifier is useful for high-confidence flagging but misses 70\% of post-2021 agents. We extract a reusable three-step leakage-audit toolkit that other blockchain-ML researchers can apply. The \texttt{defi\_hf\_trader} ablation (removing 37\% of agents labeled by behavioral criteria) shows AUC is unchanged, confirming the classifier learns from provenance-pure signal. All code, features, labels, and diagnostic scripts are released.
\end{abstract}

\keywords{on-chain AI agents, behavioral fingerprinting, label leakage, Ethereum, graph neural networks, security audit, honest baselines, reproducibility}
\maketitle

\section{1. Introduction}

\subsection{1.1 Motivation}

Autonomous AI agents are proliferating on public blockchains at unprecedented speed. Simple trading bots, complex multi-step DeFi strategy executors, MEV searchers, and LLM-driven market makers collectively control tens of billions of dollars on Ethereum mainnet alone. Platforms such as Autonolas, Fetch.ai, AI Arena, AI16Z ELIZA, and Virtuals Protocol make on-chain agent deployment increasingly turnkey. Recent reports attribute more than $45M in 2025-2026 losses to AI-agent-specific failure modes, including approval exploits, oracle manipulation, and cascading liquidation failures triggered by coordinated rebalancing.

Yet from the perspective of the Ethereum protocol, an AI-controlled externally owned account (EOA) is indistinguishable from a human-operated one. No on-chain type system exposes whether a transaction's sender is a human, a smart-contract wallet, or an AI policy running in a data center. This opacity creates three fundamental problems:

1. \textbf{Security auditors} cannot isolate agent-specific risk (runaway approvals, unbounded slippage tolerance, mis-specified LLM tool calls) from human behavior.
2. \textbf{Protocol designers} cannot offer agent-aware primitives (rate limiting, bounded approvals, composability guards) because they cannot tell who is using the protocol.
3. \textbf{Regulators and monitoring services} have no visibility into what fraction of market activity is automated, which limits their ability to reason about systemic risk.

Closing these gaps requires a reliable way to identify agents from public on-chain data alone, without relying on platform registries or voluntary disclosures. Figure~\ref{fig:pipeline} shows where this paper fits within our broader research program.

\input{../../shared/figures/fig\_pipeline.tex}

\textbf{Key Counter-Intuitive Findings.} Three results in this paper overturn conventional assumptions and motivate our methodological contributions:

\begin{enumerate}
    \item \textbf{Security audit reversal: agents are NOT riskier than humans under honest labels.} The widely cited claim that AI agents hold 10x more unlimited token approvals than humans is an artifact of heuristic labeling bias. Under provenance-verified labels, the direction reverses: human DeFi users hold more unlimited approvals (66.1 mean) than curated MEV bots (49.1 mean), because the heuristic systematically misclassified heavy DeFi users as agents (Section~4.8).
    \item \textbf{The task is harder than reported: honest temporal AUC is 0.67, not 0.98.} Our initial pipeline reported GBM AUC 0.98 on heuristically labeled data. After removing label leakage and evaluating under a temporal holdout that simulates forward deployment, the realistic AUC drops to 0.67---a 31-point gap that reveals both labeling bias and temporal distribution shift (Sections~4.1, 4.7).
    \item \textbf{The leakage-audit toolkit is a standalone contribution.} We formalize a reusable three-step audit---label--feature overlap, AUC across label-purity tiers, and cross-label-scheme transfer---that detected the leakage in our own pipeline and that other blockchain-ML researchers can apply to any behavioral classifier (Section~5.4, Algorithm~\ref{alg:leakage\_audit}).
\end{enumerate}

\subsection{1.2 The Seduction of High Numbers}

When we began this study, we built what seemed to be a compelling pipeline: 3,316 Ethereum addresses (2,590 labeled AGENT, 726 labeled HUMAN) mined from four agent platforms (Autonolas, Fetch.ai, AI Arena, Virtuals), a curated MEV-bot list, and a set of ENS-verified humans. Labels were assigned via a 4-gate heuristic we called \textbf{C1-C4}:

- C1: public registry or known identity;
- C2: behavioral consistency with the declared class;
- C3: multi-source cross-validation, operationalized as thresholds on \texttt{hour\_entropy > 2.5}, \texttt{burst\_ratio > 0.1}, and \texttt{tx\_interval\_cv < 1.0};
- C4: temporal stability across windows.

Gradient Boosting on all 23 features achieved \textbf{AUC 0.9803} (5-fold 5-rep CV, n=3316). Random Forest reached 0.9766, and even logistic regression landed at 0.9521. A Graph Attention Network on the induced transaction graph obtained 0.934, and GraphSAGE 0.967. The security audit appeared to confirm the hypothesis that agents are riskier than humans: 10.3x more unlimited approvals on average, 1.7x higher DEX interaction rate, and higher revert rates.

This is the number we would have written up in the abstract, and it is wrong as a claim about agent identification.

\subsection{1.3 The Leakage Discovery}

In the final week of experiments, we noticed that three of the top-four classifier features were algebraically identical to three C1-C4 gates:

\begin{verbatim}
hour\_entropy (C3 gate)     == active\_hour\_entropy (classifier feature)
burst\_ratio  (C3 gate)     == burst\_frequency     (classifier feature)
tx\_interval\_cv (C3 gate)   ~  tx\_interval\_std / tx\_interval\_mean (classifier features)
\end{verbatim}

In other words, the same quantities that had decided whether an address was labeled AGENT were then fed to the classifier as predictors. The resulting AUC is therefore not a measure of how well behavioral features separate agents from humans; it is a measure of how well the classifier recovers the thresholding rule the mining script used. The agents-have-10x-more-approvals headline comes from the same labeling bias: C1-C4 preferentially accepted token holders who were already active DeFi users.

Fortunately, we identified a strict set of 64 addresses with provenance labels that did not pass through C1-C4: 33 agents from the original manual pilot set (jaredfromsubway, Wintermute, 1inch resolver, Flashbots builders, Autonolas registry entries, etc.) and 31 humans from the curated ENS list (vitalik.eth, Hayden Adams, and comparable accounts). We then launched two successive expanded provenance mining campaigns (v3, v4) that grew the honest set to \textbf{1,147 addresses (533 agents, 614 humans)} by adding Gelato executor contracts, Keep3r network executors, high-frequency DeFi traders with >50 daily Uniswap calls, additional curated MEV bots, Compound V3 liquidators, Chainlink keepers, ENS-interaction-verified humans, exchange depositors, Gitcoin donors, ENS reverse setters, PoolTogether depositors, and governance voters. On this \textbf{full provenance set (n=1,147)}:

- Random Forest 10-fold AUC is \textbf{0.803};
- Gradient Boosting 10-fold AUC is \textbf{0.810};
- GAT 5-fold AUC is \textbf{0.832 +/- 0.021};
- The C1-C4 diagnostic heuristic agrees with provenance labels only \textbf{48.8\%} of the time, down from 92.9\% on the original 14-address pilot, confirming that heuristic labeling increasingly disagrees with provenance labels as the dataset grows and diversifies.

On the \textbf{strict provenance subset (n=64)}, used as a sensitivity check: RF LOO-CV AUC is 0.7713, GAT 5-fold AUC is 0.883, and GBM is indistinguishable from a majority classifier (McNemar p=0.58). The security audit direction \textbf{reverses}: famous ENS humans hold more unlimited approvals (66.1 mean) than curated MEV bots (49.1 mean) because the ENS accounts are heavier DeFi users than a typical searcher.

We therefore rewrote the paper around this discovery. This paper is not an attempt to hide the leakage; it is an attempt to turn the leakage into a methodological contribution, to present the honest baselines, and to provide a reusable audit template so that future work on behavioral classification of on-chain actors does not repeat the same mistake.

\subsection{1.4 Research Questions}

- \textbf{RQ1 (Identifiability):} Can AI agents be reliably distinguished from human users on the basis of public on-chain behavior alone, using provenance-only ground truth?
- \textbf{RQ2 (Leakage):} What is the magnitude of the gap between a heuristically labeled classifier and a provenance-only classifier, and how should future work design around it?
- \textbf{RQ3 (Model class):} Do graph neural networks, which can exploit transfer-graph structure beyond node features, beat tabular classifiers on the honest set?
- \textbf{RQ4 (Generalization):} Do classifiers trained on one platform (e.g., Autonolas) transfer to agents from a curated external set?
- \textbf{RQ5 (Security posture):} Does the "agents are riskier than humans" hypothesis survive the honest relabeling?
- \textbf{RQ6 (Temporal robustness):} How much does random CV overestimate real-world performance compared to a temporal holdout that simulates forward deployment?

\subsection{1.5 Contributions}

1. \textbf{An honest empirical baseline on 1,147 provenance-verified addresses.} Random Forest 10-fold AUC = 0.803, GBM 10-fold AUC = 0.810, and Graph Attention Network 5-fold AUC = 0.832 on 1,147 provenance-verified Ethereum addresses (533 agents, 614 humans) drawn from 11 agent sources and 7 human sources. A strict 64-address subset (33 agents, 31 humans) serves as a sensitivity check, yielding RF LOO-CV AUC 0.7713 and GAT 0.883. These are the numbers we recommend as the reference point for future work.
2. \textbf{A leakage case study.} We document how a seemingly sensible 4-gate labeling oracle injected three top-4 classifier features into the labels, causing an 18-point AUC inflation (0.9803 to 0.803). We release the diagnostic scripts as a template (\texttt{verify\_c1c4.py}, \texttt{pipeline\_provenance.py}, \texttt{statistical\_tests.py}, \texttt{mine\_addresses\_v4.py}).
3. \textbf{Statistical significance on the honest set.} McNemar, DeLong, and bootstrap tests on the 64-address strict subset show that Random Forest significantly beats Gradient Boosting (McNemar p=0.0225, DeLong p=0.0014, bootstrap CI95 [-0.24, -0.06]), GBM is indistinguishable from a majority baseline, and a single best feature (tx\_interval\_mean) matches or beats GBM.
4. \textbf{Cross-platform evaluation.} Autonolas-to-trusted transfer AUC is 0.24 to 0.34, worse than chance, while Fetch.ai-to-AI-Arena transfer AUC is 0.98, exposing that the apparent high performance comes from intra-platform correlations rather than a generalizable agent concept.
5. \textbf{GNN vs tabular comparison.} On the full provenance set (n=1,147), GAT (AUC 0.832) beats all tabular classifiers. On the strict subset (n=64), GAT reaches 0.883, beating GraphSAGE (0.784) and all tabular models, suggesting that graph structure carries signal beyond per-node features.
6. \textbf{Combined feature study.} Extending the 23 Paper 1 features with 8 AI-specific features from our companion work (Paper 3) yields a combined LightGBM 5x3-CV AUC of 0.7172 on the trusted subset, essentially no improvement (delta +0.013) over Paper 1 features alone.
7. \textbf{4-dimensional security audit with honest reversal.} We report both the leaky-set (agents 10.3x more unlimited approvals) and trusted-set (humans hold more) audits side by side, and explain why the direction reverses.
8. \textbf{C1-C4 diagnostic divergence.} The C1-C4 heuristic agrees with provenance labels only 48.8\% on n=1,147 (down from 55.8\% on n=563 and 92.9\% on the 14-address v2 pilot), demonstrating that heuristic labeling increasingly disagrees with ground truth as the dataset diversifies.
9. \textbf{Temporal holdout evaluation.} A temporal split at median first-seen block (April 2021) shows random CV overestimates AUC by 13--16 points; precision remains high (0.81) but recall collapses to 0.31 on post-2021 agents, quantifying the deployment gap.
10. \textbf{A reusable three-step leakage-audit toolkit (Algorithm~\ref{alg:leakage\_audit}).} We formalize the audit into: (1) label--feature overlap via Jaccard/correlation, (2) AUC comparison across label-purity tiers, and (3) cross-label-scheme transfer. This is a standalone methodological contribution that other blockchain-ML researchers can cite and apply.
11. \textbf{Honest baseline contextualization.} We compare our provenance-verified AUC (0.803) with published bot-detection numbers across platforms (Twitter 0.75--0.85, Reddit $\sim$0.80, Ethereum Ponzi 0.92 with leaky labels), showing that honest on-chain agent identification is competitive in the clean-label regime.
12. \textbf{Open artifacts.} Code, mined features, labels, statistical test scripts, GNN model weights, and the full leaky-vs-honest comparison JSONs are released.

\subsection{1.6 Paper Roadmap}

Section 2 surveys related work. Section 3 describes the mining pipeline, the 23 features, the tabular classifiers, the GNN architectures, and the security audit. Section 4 presents nine result subsections: honest tabular performance (4.1), statistical significance (4.2), cross-platform generalization (4.3), GNN vs tabular (4.4), combined Paper 1 + Paper 3 features (4.5), SHAP explanations (4.6), temporal holdout evaluation (4.7), the 4-dimensional security audit (4.8), and a comparison with published bot-detection baselines (4.9). Section 5 is the threats-to-validity section and is the methodological heart of the paper; it walks through the C1-C4 leakage discovery and presents a reusable three-step leakage-audit toolkit as Algorithm~\ref{alg:leakage\_audit}. Section 6 discusses practical implications for auditors and regulators, limitations, and a clean re-mining roadmap. Section 7 concludes.


\section{2. Related Work}

Our study bridges four research areas: bot detection on blockchains, MEV and automated trading, AI agent platforms, and blockchain security analysis. We also draw on the broader machine-learning literature on label leakage, data-snooping bias, and cross-domain evaluation.

\subsection{2.1 Bot Detection on Blockchains}

Detecting automated activity on blockchains has been studied primarily in the context of specific malicious behaviors. Chen et al. [1] used transaction features and opcode analysis to detect Ponzi schemes on Ethereum, demonstrating that behavioral patterns in transaction data reveal automated activity. Victor and Weintraud [2] detected wash trading on decentralized exchanges via graph features. Torres et al. [3] studied Ethereum frontrunner bots through mempool analysis, identifying displacement, insertion, and suppression attacks. Victor et al. [4] proposed behavioral fingerprinting for blockchain bots using temporal and gas features similar in spirit to our temporal and gas groups, but restricted to MEV bots. Li et al. [5] studied Telegram trading bots and their on-chain footprint.

\textbf{Gap.} Existing methods target specific malicious behaviors (wash trading, Ponzi schemes, frontrunning) or specific bot families (MEV, Telegram), not the general class of AI-controlled addresses. More importantly, none of these works explicitly audit the relationship between their labeling oracle and their classifier features; our leakage discovery suggests such audits are missing from the field.

\subsection{2.2 MEV and Automated Trading}

Daian et al. [6] introduced MEV in "Flash Boys 2.0" and demonstrated that transaction ordering creates a searcher-bot ecosystem. Weintraub et al. [7] measured Flashbots MEV at scale. Qin et al. [8] quantified MEV across DeFi protocols, showing sandwich attacks, liquidations, and arbitrage extract billions annually. Park et al. [9] analyzed post-Merge MEV and the effect of Proposer-Builder Separation. Gupta et al. [10] studied cross-domain MEV across chains and L2s.

\textbf{Gap.} MEV research quantifies the economic impact of automation but does not attempt general-purpose identification. Our classifier is meant to generalize beyond MEV searchers to DeFi management agents, portfolio rebalancers, and LLM-driven strategy agents.

\subsection{2.3 AI Agent Platforms}

Autonolas (OLAS) [11] provides an open-source framework for autonomous agent services, with on-chain registration. AI16Z ELIZA [12] enables AI agents to interact with blockchain protocols through natural language. Virtuals Protocol [13] offers a launchpad for tokenized AI agents. Fetch.ai [14], SingularityNET [15], and AI Arena [16] contribute to the growing on-chain agent population. These platforms deploy agents but do not provide cross-platform identification. A recent survey by He et al. [23] catalogs the taxonomy of on-chain AI agents but provides no empirical identification methodology.

\subsection{2.4 Blockchain Security Analysis}

Zhou et al. [17] systematized DeFi attacks; Wen et al. [18] analyzed DeFi composability risk; Grech et al. [19] and Tsankov et al. [20] built static analysis tools for smart contracts. Wang et al. [21] studied the risks of AI agents with cryptocurrency wallet access, identifying prompt injection as a primary attack vector. Zhang et al. [22] analyzed MaxUint256 approval attacks.

\textbf{Gap.} Existing security analysis focuses on protocol-level vulnerabilities. No prior work assesses the security posture of AI agent addresses as a class or compares it to a human baseline, and the handful of works that attempt to do so do not control for labeling bias.

\subsection{2.5 Label Leakage and Data-Snooping Bias}

Kaufman et al. [24] formalized leakage in predictive modeling and catalog several subtle forms. Roelofs et al. [25] studied test-set reuse as an implicit form of leakage in ML benchmarks. Geirhos et al. [26] documented shortcut learning: deep networks exploit unintended features that correlate with the label during training but not at deployment. In the blockchain space, Tikhomirov and Ivanitskiy [27] warned about address-type features becoming labels in phishing classifiers. Our work extends this line: we show that a seemingly principled 4-gate labeling heuristic can encode three classifier features into the label, producing a 25-point AUC inflation that is invisible to standard cross-validation.

\subsection{2.6 Graph Neural Networks on Transaction Graphs}

GraphSAGE [28] and GAT [29] are standard inductive GNN architectures. In the blockchain setting, Weber et al. [30] applied GNNs to Bitcoin anti-money-laundering. Patel et al. [31] used GNNs to detect Ethereum phishing. Shamsi et al. [32] surveyed GNNs for blockchain analytics. Most of these works report aggregate accuracy on full labeled graphs; very few report transfer performance across labeling schemes, and none to our knowledge report the honest-vs-leaky gap that is the central finding of this paper.

\subsection{2.7 Our Position}

Relative to the above, this paper offers: (1) a general behavioral identification framework rather than a specific bot-class detector; (2) an honest provenance-only baseline suitable for comparison across future works; (3) a documented leakage case study that other blockchain-ML researchers can use as a negative-result template; (4) a direct tabular-vs-GNN comparison on the same ground truth; (5) the first public honest security audit of AI-agent-labeled Ethereum addresses.


\section{3. Methodology}

\subsection{3.1 Data Collection}

We collect transaction data from Ethereum mainnet via Etherscan V2 API using a 6-key rotation strategy that achieves approximately 30 calls/sec sustained throughput. For each address we extract up to 5,000 most recent transactions covering normal transactions, internal transactions, ERC-20 transfers, and approval events. All data is public.

\subsection{3.2 Address Mining and Labeling}

\subsubsection{3.2.1 Strategy A: Platform Holders (C1-C4 labeled, \textbf{used for inference only})}

We identify addresses holding tokens of four agent platforms (Autonolas OLAS, Fetch.ai FET, AI Arena NRN, Virtuals VIRT). This yields \textbf{3,252 platform-token-holder addresses}. Each holder is labeled via our original C1-C4 gate:

- \textbf{C1 (public registry):} the address appears in a public platform registry or has a known identity.
- \textbf{C2 (behavioral consistency):} the address's on-chain behavior matches the declared class.
- \textbf{C3 (multi-source cross-validation):} operationalized as thresholds on \texttt{hour\_entropy > 2.5}, \texttt{burst\_ratio > 0.1}, \texttt{tx\_interval\_cv < 1.0}.
- \textbf{C4 (temporal stability):} the behavioral class is stable across time windows.

\subsubsection{3.2.2 Strategy B: Strict Provenance Subset (n=64, \textbf{sensitivity check})}

We identify 64 addresses whose labels do not depend on any behavioral threshold:

- \textbf{33 AGENT addresses} from (a) the original Paper 1 pilot set of manually curated bots (jaredfromsubway, Wintermute, 1inch resolver, Flashbots builder, Autonolas service operators), (b) the Paper 0 Delphi-validated agent set, and (c) an external curated MEV-bot list.
- \textbf{31 HUMAN addresses} from the ENS-verified human list (vitalik.eth, Hayden Adams, nick.eth, and comparable accounts whose identity is publicly verifiable).

None of the trusted labels pass through \texttt{hour\_entropy}, \texttt{burst\_ratio}, or \texttt{tx\_interval\_cv} thresholds. This subset serves as a strict sensitivity check.

\subsubsection{3.2.3 Strategy C: Expanded Provenance Mining v4 (n=1,147, \textbf{primary honest set})}

To address the sample-size limitation of n=64, we launched two successive expanded provenance mining campaigns (v3, v4) that grew the honest set to \textbf{1,147 addresses (533 agents, 614 humans)} across 18 source categories, all verified through non-behavioral provenance. The v4 campaign (\texttt{mine\_addresses\_v4.py}) added five new provenance sources: Compound V3 liquidators (97), Chainlink keepers (108), Gitcoin donors (146), ENS reverse setters (144), and PoolTogether depositors (89).

\textbf{Agent sources (533 addresses, 11 categories):}

\begin{table}[h]
\centering
\small
\begin{tabular}{lll}
\toprule
Source category & Count & Provenance basis \\
\midrule
defi\_hf\_trader & 199 & Addresses with >50 daily Uniswap router calls, verified via on-chain call frequency \\
chainlink\_keeper & 108 & Chainlink Automation / Keeper registry-attested upkeep performers \\
compound\_v3\_liquidator & 97 & Compound V3 liquidation callers, verified via Comet contract events \\
gelato\_executor & 47 & Gelato Network executor contracts, verified via Gelato registry \\
keep3r\_executor & 46 & Keep3r Network keeper addresses, verified via Keep3r bond registry \\
pilot\_agent & 15 & Original manually curated pilot set \\
curated\_mev\_bot & 10 & External curated MEV-bot list with known operator identities \\
expanded\_mev\_bot & 10 & Additional MEV bots verified via Flashbots builder registry \\
flash\_loan\_user & 1 & Address verified as Aave/dYdX flash loan initiator \\
\bottomrule
\end{tabular}
\end{table}

\textbf{Human sources (614 addresses, 7 categories):}

\begin{table}[h]
\centering
\small
\begin{tabular}{lll}
\toprule
Source category & Count & Provenance basis \\
\midrule
human\_ens\_interaction & 167 & Addresses that registered or renewed ENS names, verified via ENS registry events \\
gitcoin\_donor & 146 & Gitcoin Grants round donors, verified via Gitcoin contract events \\
ens\_reverse\_setter & 144 & Addresses that set ENS reverse records, verified via ENS reverse registrar \\
pooltogether\_depositor & 89 & PoolTogether prize pool depositors, verified via deposit events \\
human\_exchange\_depositor & 28 & Addresses depositing to known exchange hot wallets, verified via deposit patterns \\
expanded\_human & 25 & Additional ENS-verified humans matched on DeFi exposure \\
curated\_human & 7 & Original curated human list (vitalik.eth, etc.) \\
human\_ens\_interaction\_governance\_voter & 5 & ENS holders who also voted in on-chain governance (Compound, Uniswap, etc.) \\
pilot\_human & 3 & Original pilot human set \\
\bottomrule
\end{tabular}
\end{table}

The v4 mining applies 23 features to all 1,147 addresses. The agent-to-human ratio is 0.87:1 (533:614), approximately balanced. All 64 addresses from Strategy B are included in Strategy C.

\textbf{Label quality note on \texttt{defi\_hf\_trader}.} The largest agent category (199 addresses, 17\% of the total) is labeled "agent" because the address makes >50 Uniswap calls per active day. This is a behavioral criterion, not a provenance-pure label: it is possible that some of these addresses are humans using aggressive trading interfaces rather than autonomous agents. We discuss this threat in Section 5.

\subsubsection{3.2.4 Exclusion Rules}

We exclude contract addresses (Uniswap routers, 0x Exchange Proxy, Seaport, Blur), exchange hot wallets (MEXC, Binance 14/16), and any address whose label cannot be pinned to a non-behavioral source.

\subsubsection{3.2.5 Decision: Demote C1-C4 to Evaluation}

After discovering that C1-C4 gates overlap with three top classifier features (see Section 5), we demote the 3,252 Strategy A addresses from training to inference-only. The 1,147 full provenance addresses (Strategy C) become the \textbf{primary honest train+test set}, with the 64 strict provenance addresses (Strategy B) as a sensitivity check. We continue to report leaky-set numbers for transparency and for statistical comparison.

The resulting three-tier evaluation hierarchy is:

\begin{table}[h]
\centering
\small
\begin{tabular}{lllll}
\toprule
Tier & Label & n & Label source & Role \\
\midrule
1 (leaky) & C1-C4 heuristic & 3,316 & Behavioral thresholds & Inference-only, not cited as headline \\
2 (full provenance) & Strategy C v4 & 1,147 & Non-behavioral provenance & \textbf{Primary honest set} \\
3 (strict provenance) & Strategy B & 64 & Manual curation & Sensitivity check \\
\bottomrule
\end{tabular}
\end{table}

\subsection{3.3 Feature Engineering: 23 Behavioral Features}

We compute 23 features organized into four groups. Let $T = \{t_1, \ldots, t_n\}$ be the transaction timestamps, $G = \{g_1, \ldots, g_n\}$ the gas prices, and $C = \{c_1, \ldots, c_n\}$ the target contracts.

\textbf{Group 1: Temporal Features (7).}

1. \texttt{tx\_interval\_mean}: mean of consecutive inter-transaction intervals.
2. \texttt{tx\_interval\_std}: standard deviation of intervals.
3. \texttt{tx\_interval\_skewness}: skewness of the interval distribution.
4. \texttt{active\_hour\_entropy}: Shannon entropy of transactions bucketed into 24 UTC hours.
5. \texttt{night\_activity\_ratio}: fraction of transactions in UTC 0:00-6:00.
6. \texttt{weekend\_ratio}: fraction of transactions on Saturday/Sunday.
7. \texttt{burst\_frequency}: fraction of consecutive transaction pairs with interval < 10 s.

\textbf{Group 2: Gas Behavior Features (6).}

8. \texttt{gas\_price\_round\_number\_ratio}: fraction of transactions with integer-Gwei gas prices.
9. \texttt{gas\_price\_trailing\_zeros\_mean}: mean number of trailing zeros in the decimal gas-price representation.
10. \texttt{gas\_limit\_precision}: gas\_limit relative to gas\_used.
11. \texttt{gas\_price\_cv}: coefficient of variation of gas prices.
12. \texttt{eip1559\_priority\_fee\_precision}: precision of priority fees in EIP-1559 transactions.
13. \texttt{gas\_price\_nonce\_correlation}: Pearson correlation between gas price and nonce.

\textbf{Group 3: Interaction Pattern Features (5).}

14. \texttt{unique\_contracts\_ratio}: unique target contracts divided by total transactions.
15. \texttt{top\_contract\_concentration}: Herfindahl-Hirschman index of interaction targets.
16. \texttt{method\_id\_diversity}: unique function signatures per contract call.
17. \texttt{contract\_to\_eoa\_ratio}: fraction of transactions whose target is a contract.
18. \texttt{sequential\_pattern\_score}: n-gram repetition score of the action sequence.

\textbf{Group 4: Approval \& Security Features (5).}

19. \texttt{unlimited\_approve\_ratio}: fraction of approve() calls that set MaxUint256.
20. \texttt{approve\_revoke\_ratio}: revoke operations per approve operation.
21. \texttt{unverified\_contract\_approve\_ratio}: fraction of approvals to Etherscan-unverified contracts.
22. \texttt{multi\_protocol\_interaction\_count}: count of distinct DeFi protocols interacted with.
23. \texttt{flash\_loan\_usage}: flash-loan indicator.

\textbf{Critical note on leakage.} Features 1-7 (temporal) and especially 2, 4, 7 are algebraically entangled with the C3 gate that labeled the Strategy A subset. This is the central issue analyzed in Section 5.

\subsection{3.4 Tabular Classifiers}

We train four classifiers via scikit-learn and LightGBM:

- \textbf{GradientBoosting (GBM):} 200 trees, max\_depth=3, lr=0.1.
- \textbf{RandomForest (RF):} 100 trees, unlimited depth.
- \textbf{LogisticRegression (LR):} L2-regularized, liblinear solver.
- \textbf{LightGBM:} 500 estimators, num\_leaves=31.

Features are standardized before training. We evaluate with two complementary schemes:

- \textbf{Leave-One-Out Cross-Validation (LOO-CV):} 64 iterations, one held-out sample per fold.
- \textbf{Repeated Stratified 5-Fold CV (5x10-CV):} 50 evaluations, mean and std reported.

Metrics: AUC-ROC (primary), precision, recall, F1, accuracy, and Brier score for calibration.

\subsection{3.5 Graph Neural Networks}

We construct a directed transaction graph $G = (V, E)$ on the 3,316-address set, with edges from sender to receiver for any transfer of value or token. Node features are the 23 behavioral features above. The resulting graph has 3,316 nodes, 2,387 edges, and is moderately sparse. For the full provenance evaluation (n=1,147), we also construct a subgraph with 1,147 nodes and 905 edges.

We train two GNN architectures:

- \textbf{GraphSAGE:} two-layer, hidden dim 64, mean aggregator, dropout 0.3, Adam lr=0.01, early stopping on validation loss.
- \textbf{GAT (Graph Attention Network):} two-layer, 8 attention heads, hidden dim 64, dropout 0.3.

We evaluate on three splits: \texttt{full\_3316} (training labels are C1-C4 leaky), \texttt{provenance\_1147} (training labels are full provenance v4; the primary honest evaluation), and \texttt{trusted\_64} (training labels are strict provenance-only; sensitivity check).

\subsection{3.6 Baselines and Ablations}

- \textbf{Majority classifier:} predicts the majority class (AGENT).
- \textbf{Single best feature:} logistic classifier on \texttt{tx\_interval\_mean}, the top univariate AUC feature on the trusted set.
- \textbf{Heuristic rule:} \texttt{burst\_frequency > 0.1 => AGENT}.
- \textbf{Feature-group ablations:} retain one group (temporal / gas / interaction / approval) at a time.

\subsection{3.7 Statistical Tests}

On the 64-address trusted set, we compute:

- \textbf{McNemar's test} for paired predictions.
- \textbf{DeLong's test} for paired AUC comparison with bootstrap variance.
- \textbf{10,000-sample bootstrap} for AUC difference confidence intervals.
- \textbf{Precision @ K} for K in {5, 10, 20, 30}.
- \textbf{Brier score} for calibration.

\subsection{3.8 Cross-Platform Evaluation}

We run four transfer splits:

1. \texttt{auto\_train\_to\_trusted\_test}: train on 1,780 Autonolas holders, test on 41 trusted addresses.
2. \texttt{trusted\_train\_to\_auto\_test}: train on 41 trusted addresses, test on 1,780 Autonolas holders.
3. \texttt{fetch\_to\_ai\_arena}: train on 923 Fetch.ai, test on 549 AI Arena.
4. \texttt{ai\_arena\_to\_fetch}: train on 549 AI Arena, test on 923 Fetch.ai.

Splits 1-2 are cross-label-scheme (provenance vs C1-C4). Splits 3-4 are intra-label-scheme (both C1-C4). The gap between these two families is a direct measurement of labeling-scheme transfer.

\subsection{3.9 Combined Paper 1 + Paper 3 Features}

To evaluate whether our 23 behavioral features can be meaningfully extended, we compute 8 AI-specific features from the companion Paper 3 work (\texttt{hour\_entropy}, \texttt{behavioral\_consistency}, \texttt{response\_latency\_variance}, \texttt{action\_sequence\_perplexity}, \texttt{gas\_nonce\_gap\_regularity}, \texttt{eip1559\_tip\_precision}, \texttt{gas\_price\_precision}, \texttt{error\_recovery\_pattern}). We train four classifiers on three feature sets: Paper 1 only (23), Paper 3 only (8), and combined (31).

\subsection{3.10 4-Dimensional Security Audit}

For each audited address we compute:

- \textbf{Dimension 1: Permission exposure.} Counts of active approvals, unlimited (MaxUint256) approvals, and unverified approval targets; average and max approval ages in days.
- \textbf{Dimension 2: Agent network topology.} Degree centrality, clustering coefficient, community membership in the transfer graph.
- \textbf{Dimension 3: MEV / DEX exposure.} DEX interaction rate, swap rate.
- \textbf{Dimension 4: Failure and gas analysis.} Revert rate, gas-used efficiency.

We run the audit twice: once on the full 3,302-address set (14 dropped due to empty histories) and once on the 50-address trusted subset. The difference between the two is a direct measurement of how labeling bias corrupts security conclusions.


\section{4. Results}

We now present the honest results in six subsections. Throughout this section, \textbf{the full provenance set (n=1,147) is the headline}. The strict provenance subset (n=64) is reported as a sensitivity check. The full n=3,316 numbers are reported only as a leaky reference point.

\subsection{4.1 Honest Performance on Provenance-Verified Sets}

\subsubsection{4.1.1 Primary: Full Provenance Set (n=1,147, 533 agents, 614 humans)}

Table 1 reports the core honest performance on the full provenance v4 set.

\textbf{Table 1. Tabular and GNN classifier performance on the full provenance set (n=1,147, 533 agents, 614 humans).}

\begin{table}[h]
\centering
\small
\begin{tabular}{lllllll}
\toprule
Model & Split & AUC-ROC & Precision & Recall & F1 & Accuracy \\
\midrule
RandomForest & 10-fold & \textbf{0.8028} & 0.7444 & 0.6229 & 0.6782 & 0.7254 \\
RandomForest & 5x10-CV & 0.8029 +/- 0.027 & -- & -- & -- & -- \\
GradientBoosting & 10-fold & 0.8095 & -- & -- & 0.7030 & 0.7341 \\
GAT & 5-fold & \textbf{0.8318 +/- 0.021} & -- & -- & 0.674 & -- \\
\bottomrule
\end{tabular}
\end{table}

\textbf{Top 10 RF features by importance (n=1,147):}

\begin{table}[h]
\centering
\small
\begin{tabular}{lll}
\toprule
Rank & Feature & Importance \\
\midrule
1 & gas\_price\_trailing\_zeros\_mean & 0.1448 \\
2 & gas\_price\_round\_number\_ratio & 0.1302 \\
3 & unique\_contracts\_ratio & 0.0927 \\
4 & top\_contract\_concentration & 0.0870 \\
5 & method\_id\_diversity & 0.0725 \\
6 & tx\_interval\_mean & 0.0675 \\
7 & gas\_price\_cv & 0.0650 \\
8 & sequential\_pattern\_score & 0.0484 \\
9 & contract\_to\_eoa\_ratio & 0.0435 \\
10 & tx\_interval\_std & 0.0345 \\
\bottomrule
\end{tabular}
\end{table}

Note that gas features now dominate the importance ranking on n=1,147, whereas temporal features (especially \texttt{active\_hour\_entropy}) dominated on the leaky set. This shift confirms that the C1-C4 leakage inflated the apparent importance of temporal features.

\textbf{RF confusion matrix (n=1,147, 10-fold aggregate):}

\begin{table}[h]
\centering
\small
\begin{tabular}{lll}
\toprule
 & Pred Human & Pred Agent \\
\midrule
True Human & 500 & 114 \\
True Agent & 201 & 332 \\
\bottomrule
\end{tabular}
\end{table}

The classifier has moderate agent recall (332/533 = 62.3\%) and high human precision (500/614 = 81.4\%), reflecting the approximately balanced class distribution (0.87:1 agent-to-human ratio).

\textbf{Feature-group ablation (RF 10-fold on n=1,147):}

\begin{table}[h]
\centering
\small
\begin{tabular}{lllll}
\toprule
Feature group & n\_feat & AUC (alone) & AUC (without) & Delta \\
\midrule
interaction & 5 & \textbf{0.751} & 0.748 & +0.047 \\
gas & 6 & 0.713 & 0.767 & +0.028 \\
approval\_security & 5 & 0.656 & 0.789 & +0.005 \\
temporal & 7 & 0.623 & 0.799 & -0.005 \\
\textbf{all 23} & 23 & \textbf{0.794} & — & — \\
\bottomrule
\end{tabular}
\end{table}

Interaction features (unique\_contracts\_ratio, top\_contract\_concentration) are the most important group on the honest set, while temporal features contribute least — removing them actually helps by +0.005. This reverses the leaky-set finding where temporal features dominated via \texttt{active\_hour\_entropy}. The reversal is further evidence that the C1-C4 leakage was concentrated in temporal features.

\textbf{C1-C4 diagnostic agreement.} We applied the C1-C4 heuristic to all 1,147 provenance-labeled addresses. The agreement rate is \textbf{48.8\%}, below chance. This is a further drop from 55.8\% on the v3 set (n=563) and 92.9\% on the original 14-address v2 pilot, demonstrating that the C1-C4 heuristic's reliability degrades as the dataset grows and diversifies. As the dataset expands to include Compound V3 liquidators, Chainlink keepers, Gitcoin donors, and PoolTogether depositors, the heuristic's behavioral thresholds increasingly mislabel addresses.

\subsubsection{4.1.2 Sensitivity Check: Strict Provenance Subset (n=64, 33 agents, 31 humans)}

Table 2 reports the strict provenance subset results under LOO-CV and repeated 5x10-CV.

\textbf{Table 2. Tabular classifier performance on the strict provenance subset (n=64, 33 agents, 31 humans).}

\begin{table}[h]
\centering
\small
\begin{tabular}{lllllll}
\toprule
Model & Split & AUC-ROC & Precision & Recall & F1 & Accuracy \\
\midrule
RandomForest & LOO & \textbf{0.7713} & 0.7419 & 0.6970 & 0.7188 & 0.7188 \\
RandomForest & 5x10-CV & \textbf{0.8030} +/- 0.102 & 0.8008 & 0.7138 & 0.7363 & 0.7447 \\
GradientBoosting & LOO & 0.6276 & 0.5882 & 0.6061 & 0.5970 & 0.5781 \\
GradientBoosting & 5x10-CV & 0.6940 +/- 0.119 & 0.6731 & 0.6638 & 0.6548 & 0.6499 \\
LogisticRegression & LOO & 0.6100 & 0.6667 & 0.6061 & 0.6349 & 0.6406 \\
LogisticRegression & 5x10-CV & 0.6203 +/- 0.125 & 0.6974 & 0.6238 & 0.6459 & 0.6614 \\
LightGBM & LOO & 0.7097 & -- & -- & 0.7273 & 0.7188 \\
LightGBM & 5x10-CV & 0.7040 +/- 0.088 & -- & -- & 0.6793 & 0.6765 \\
\bottomrule
\end{tabular}
\end{table}

The n=64 results are broadly consistent with n=1,147: RF is the best tabular model, GBM underperforms, and standard deviations are wide. The slightly lower RF AUC on n=64 (0.7713 LOO vs 0.803 10-fold on n=1,147) may reflect the reduced statistical power of the smaller set rather than a genuine performance difference.

\textbf{Key findings across both sets.}

1. \textbf{Random Forest is the best honest tabular model} across both set sizes (10-fold AUC 0.803 on n=1,147, LOO AUC 0.7713 on n=64).
2. \textbf{GBM underperforms RF} on both honest sets, despite being nearly identical on the leaky set. We confirm this gap with statistical tests in Section 4.2.
3. \textbf{Standard deviations shrink substantially from n=64 to n=1,147} (RF repeated-CV std drops from 0.102 to 0.027), reflecting the benefit of a larger honest corpus.
4. \textbf{Feature importance shifts.} On n=1,147, gas features (gas\_price\_trailing\_zeros\_mean, gas\_price\_round\_number\_ratio) dominate; on the leaky set, temporal features (active\_hour\_entropy) dominated. This shift is further evidence of leakage.

\textbf{Leaky reference (do not cite as headline).} On the full 3,316 addresses with C1-C4 labels, GBM achieves 5-fold 5-rep CV AUC 0.9803 +/- 0.006, RF 0.9766 +/- 0.006, LR 0.9521 +/- 0.012, LightGBM 0.9822 +/- 0.006. The gap between 0.98 and 0.80 is the leakage magnitude.

\textbf{Per-feature univariate AUC (trusted n=64).} Table 3 lists the top features by Mann-Whitney-derived AUC.

\textbf{Table 3. Top univariate AUC on the strict trusted set (n=64).}

\begin{table}[h]
\centering
\small
\begin{tabular}{llll}
\toprule
Rank & Feature & Univariate AUC & Group \\
\midrule
1 & tx\_interval\_mean & 0.7683 & Temporal \\
2 & tx\_interval\_std & 0.7527 & Temporal \\
3 & unlimited\_approve\_ratio & 0.7116 & Approval \\
4 & unverified\_contract\_approve\_ratio & 0.7107 & Approval \\
5 & gas\_price\_round\_number\_ratio & 0.7072 & Gas \\
6 & gas\_price\_trailing\_zeros\_mean & 0.7067 & Gas \\
7 & active\_hour\_entropy & 0.7028 & Temporal \\
8 & multi\_protocol\_interaction\_count & 0.6960 & Approval \\
9 & unique\_contracts\_ratio & 0.6950 & Interaction \\
10 & method\_id\_diversity & 0.6911 & Interaction \\
\bottomrule
\end{tabular}
\end{table}

Contrast with the leaky set, where \texttt{active\_hour\_entropy} led with AUC 0.9347 and \texttt{tx\_interval\_skewness} reached 0.872. The collapse of these features from 0.93 to 0.70 is direct evidence that the leaky AUCs were inflated by labeling overlap.

\subsection{4.2 Statistical Significance Tests}

Given n=64, we report three complementary tests for each pairwise model comparison: McNemar's test on paired predictions, DeLong's test for paired AUCs, and a 10,000-sample bootstrap CI for the AUC difference.

\textbf{Table 4. Pairwise statistical tests on the strict trusted set (n=64, LOO predictions).}

\begin{table}[h]
\centering
\small
\begin{tabular}{lllll}
\toprule
Comparison & McNemar p & DeLong p & Bootstrap CI95 AUC diff & Significant? \\
\midrule
GBM vs RF & \textbf{0.0225} & \textbf{0.0014} & [-0.24, -0.06] & \textbf{Yes} (RF wins) \\
GBM vs LR & 0.4545 & 0.7895 & [-0.11, 0.15] & No \\
GBM vs Single-Feature & 0.7353 & <0.001 & [-0.52, -0.24] & Mixed \\
GBM vs Majority & 0.5839 & -- & -- & \textbf{No} \\
\bottomrule
\end{tabular}
\end{table}

\textbf{Key findings.}

1. \textbf{Random Forest significantly beats Gradient Boosting} on the trusted set by three independent tests. This is a qualitative reversal of the leaky-set conclusion where GBM and RF were statistically indistinguishable.
2. \textbf{GBM is NOT significantly better than a majority classifier} (McNemar p=0.58). Read literally: a model that always predicts AGENT performs as well as the GBM on this set.
3. \textbf{A single feature (\texttt{tx\_interval\_mean}) achieves AUC 1.00 at the threshold level} but low accuracy (0.53). This is an artifact of degenerate thresholding on small n; it is reported for completeness but not as a headline claim. The more honest comparison is that RF beats every other model we tried.
4. \textbf{Brier scores:} RF 0.194, LR 0.274, LightGBM not reported, GBM 0.355. RF is the best-calibrated model as well as the most accurate.

\subsection{4.3 Cross-Platform Generalization}

Table 5 reports the four cross-platform splits. The critical contrast is between splits 1-2 (cross-label-scheme: provenance vs C1-C4) and splits 3-4 (intra-label-scheme: both C1-C4).

\textbf{Table 5. Cross-platform transfer AUC.}

\begin{table}[h]
\centering
\small
\begin{tabular}{lllllll}
\toprule
Split & Direction & GBM AUC & RF AUC & LR AUC & n\_train & n\_test \\
\midrule
Autonolas -> Trusted (honest) & C1-C4 -> provenance & 0.269 & 0.243 & 0.578 & 1,780 & 41 \\
Trusted -> Autonolas & provenance -> C1-C4 & 0.347 & 0.341 & 0.603 & 41 & 1,780 \\
Fetch.ai -> AI Arena & C1-C4 -> C1-C4 & \textbf{0.980} & 0.969 & 0.962 & 923 & 549 \\
AI Arena -> Fetch.ai & C1-C4 -> C1-C4 & \textbf{0.992} & 0.975 & 0.973 & 549 & 923 \\
\bottomrule
\end{tabular}
\end{table}

\textbf{Key findings.}

1. \textbf{Cross-label-scheme transfer collapses to worse than chance} (GBM 0.269, RF 0.243). A classifier trained on Autonolas C1-C4 labels and tested on trusted provenance labels produces \textbf{anti-predictions}: it labels agents as humans and vice versa.
2. \textbf{Intra-label-scheme transfer is near-perfect} (AUC 0.97-0.99). Fetch.ai and AI Arena holders are indistinguishable under the shared C1-C4 oracle, because the oracle is learning the oracle, not the concept.
3. \textbf{The 0.27 to 0.99 gap is another expression of leakage.} It tells us the apparent high transfer performance is spurious: it measures consistency of the labeling heuristic, not consistency of behavior across platforms.
4. \textbf{Logistic regression, paradoxically, is the best cross-label-scheme model} (0.578 / 0.603). We interpret this as follows: LR cannot fit the leaky labels as tightly as GBM/RF, so it suffers less when the labels change.

\textbf{Implication.} Any future paper claiming high agent-identification AUC on platform-holder data must include this cross-label-scheme transfer test, or the result cannot be distinguished from leakage.

\subsection{4.4 GNN vs Tabular Classifiers}

We train GraphSAGE and GAT on the 3,316-node transaction graph and evaluate on the three splits. Table 6 reports the results.

\textbf{Table 6. GNN vs tabular AUC across three evaluation tiers.}

\begin{table}[h]
\centering
\small
\begin{tabular}{llll}
\toprule
Model & full\_3316 AUC (5-fold) & provenance\_1147 AUC (5-fold) & trusted\_64 AUC (5-fold) \\
\midrule
GraphSAGE & 0.9666 +/- 0.009 & -- & 0.7841 +/- 0.084 \\
GAT & 0.9341 +/- 0.011 & \textbf{0.8318 +/- 0.021} & \textbf{0.8825 +/- 0.108} \\
RF (tabular) & 0.9766 +/- 0.006 & 0.8028 (10-fold) & 0.8030 +/- 0.102 (5x10) \\
GBM (tabular) & 0.9803 +/- 0.006 & 0.8095 (10-fold) & 0.6940 +/- 0.119 (5x10) \\
\bottomrule
\end{tabular}
\end{table}

\textbf{Key findings.}

1. \textbf{GAT is the best honest model at both scales.} On the full provenance set (n=1,147), GAT reaches 5-fold AUC 0.8318 +/- 0.021, beating RF 0.8028 and GBM 0.8095. On the strict subset (n=64), GAT reaches 0.8825 +/- 0.108 vs RF 0.8030. The higher AUC on n=64 is partly due to the smaller, more homogeneous set and partly due to high fold-to-fold variance (std 0.108 vs 0.021 on n=1,147).
2. \textbf{GAT variance is dramatically lower on n=1,147.} The std drops from 0.108 (n=64, ~13 test nodes per fold) to 0.021 (n=1,147, ~229 test nodes per fold), making the n=1,147 result far more reliable.
3. \textbf{GraphSAGE does not help on the trusted set.} GraphSAGE 0.7841 (n=64) is close to RF 0.8030 and worse than GAT. We attribute GAT's advantage to the attention mechanism, which can down-weight noisy edges in the sparse transfer graph.
4. \textbf{On the leaky full set, the ranking inverts.} GraphSAGE reaches 0.9666 and outperforms GAT 0.9341, while the tabular GBM reaches 0.9803. The leaky ranking has no predictive value for the honest setting.
5. \textbf{GNN beats tabular on all honest sets}, which is the first positive empirical result in this paper that survives the leakage correction. We consider this the honest headline: \textbf{GAT 5-fold AUC 0.832 on n=1,147} is the primary reference number.

\subsection{4.5 Combined Paper 1 + Paper 3 Features}

Table 7 reports four classifiers on three feature sets: Paper 1 only (23 behavioral features), Paper 3 only (8 AI features), and combined (31 features). The Paper 3 features are described in the companion paper.

\textbf{Table 7. Classifier AUC by feature set on the strict trusted subset (n=64, LOO and 5x3-CV).}

\begin{table}[h]
\centering
\small
\begin{tabular}{llll}
\toprule
Feature set & Model & LOO AUC & 5x3-CV mean AUC \\
\midrule
P1 only (23) & GBM & 0.6276 & 0.7267 \\
P1 only (23) & RF & 0.7713 & 0.8183 \\
P1 only (23) & LR & 0.6100 & 0.6696 \\
P1 only (23) & LightGBM & 0.7097 & 0.7040 \\
P3 only (8) & GBM & 0.7146 & 0.7008 \\
P3 only (8) & RF & 0.7292 & 0.7741 \\
P3 only (8) & LR & 0.6432 & 0.6807 \\
P3 only (8) & LightGBM & 0.7419 & 0.7392 \\
Combined (31) & GBM & 0.6393 & 0.7093 \\
Combined (31) & RF & 0.7625 & 0.8151 \\
Combined (31) & LR & 0.6256 & 0.6646 \\
Combined (31) & LightGBM & 0.6872 & \textbf{0.7172} \\
\bottomrule
\end{tabular}
\end{table}

\textbf{Key findings.}

1. \textbf{Combining the two feature sets buys almost nothing on the trusted set.} LightGBM 5x3-CV improves marginally from 0.7040 (P1-only) to 0.7172 (combined), a delta of +0.013. RF actually drops from 0.8183 to 0.8151.
2. \textbf{Paper 3 features alone are competitive.} With only 8 features, LightGBM reaches 0.7419 LOO and 0.7392 5x3-CV, close to Paper 1's 23-feature numbers. This is consistent with Paper 3's observation that \texttt{hour\_entropy} and \texttt{behavioral\_consistency} are individually strong.
3. \textbf{On the full leaky set, combined features look stronger} (GBM 0.9834 vs 0.9787 for P1 only), but this gap is within noise of the leakage artifact and should not be cited.
4. \textbf{Implication.} There is no free lunch from naively stacking feature sets on a small honest corpus. Ensemble gains require either more data or smarter late fusion.

\subsection{4.6 SHAP Feature Explanations}

We compute SHAP TreeExplainer values on the RF model trained on n=1,147 to understand feature directionality — not just importance rank but whether higher values push toward agent or human.

\textbf{Table 8a. Top 10 SHAP mean |value| (RF, n=1,147).}

\begin{table}[h]
\centering
\small
\begin{tabular}{lllll}
\toprule
Feature & Mean & SHAP &  & Direction \\
\midrule
top\_contract\_concentration & 0.057 & higher → agent &  &  \\
gas\_price\_cv & 0.044 & higher → agent &  &  \\
unique\_contracts\_ratio & 0.032 & higher → agent &  &  \\
gas\_price\_round\_number\_ratio & 0.024 & higher → agent &  &  \\
gas\_price\_trailing\_zeros\_mean & 0.023 & higher → agent &  &  \\
method\_id\_diversity & 0.017 & higher → agent &  &  \\
contract\_to\_eoa\_ratio & 0.014 & higher → agent &  &  \\
tx\_interval\_mean & 0.013 & higher → agent &  &  \\
unlimited\_approve\_ratio & 0.013 & higher → agent &  &  \\
sequential\_pattern\_score & 0.013 & higher → agent &  &  \\
\bottomrule
\end{tabular}
\end{table}

\textbf{Key finding.} \texttt{active\_hour\_entropy} — the feature that dominated the leaky set (univariate AUC 0.93) — ranks near the bottom of SHAP importance (0.002) on the honest set. This is the clearest single-feature confirmation that the C1-C4 leakage inflated temporal features. The honest model relies on gas pricing and interaction diversity, not timing patterns.

\subsection{4.7 Temporal Holdout Evaluation}

Random cross-validation allows post-2021 addresses into every training fold, leaking information about the evolving agent population into the model. To measure the true temporal generalization gap, we split the 1,147 provenance-verified addresses by median first-seen block number (block 12,321,026, approximately April 27, 2021). The \textbf{train set} consists of 573 addresses first seen before the median block (170 agents, 403 humans; 70\% human), and the \textbf{test set} consists of 574 addresses first seen after the median block (363 agents, 211 humans; 63\% agent). The class-distribution shift---from majority-human in the pre-2021 era to majority-agent in the post-2021 era---reflects the real-world proliferation of on-chain agents.

\textbf{Table 9. Random CV vs.\ temporal holdout performance on n=1,147.}

\begin{table}[h]
\centering
\small
\begin{tabular}{llllll}
\toprule
Model & Evaluation & AUC & Precision & Recall & F1 \\
\midrule
RandomForest & Random 5x10-CV & 0.803 $\pm$ 0.027 & 0.755 & 0.623 & 0.681 \\
RandomForest & Temporal holdout & 0.669 & 0.812 & 0.309 & 0.447 \\
 & $\Delta$ & \textbf{--0.134} & +0.057 & \textbf{--0.314} & --0.234 \\
\midrule
GradientBoosting & Random 5x10-CV & 0.802 $\pm$ 0.027 & 0.725 & 0.675 & 0.698 \\
GradientBoosting & Temporal holdout & 0.664 & 0.767 & 0.499 & 0.604 \\
 & $\Delta$ & \textbf{--0.137} & +0.042 & \textbf{--0.176} & --0.094 \\
\midrule
LogisticRegression & Random 5x10-CV & 0.779 $\pm$ 0.028 & 0.711 & 0.658 & 0.682 \\
LogisticRegression & Temporal holdout & 0.616 & 0.739 & 0.422 & 0.537 \\
 & $\Delta$ & \textbf{--0.163} & +0.028 & \textbf{--0.236} & --0.145 \\
\bottomrule
\end{tabular}
\end{table}

\textbf{Key findings.}

\begin{enumerate}
    \item \textbf{Random CV overestimates AUC by 13--16 points.} RF drops from 0.803 to 0.669 ($\Delta = -0.134$), GBM from 0.802 to 0.664 ($\Delta = -0.137$), and LR from 0.779 to 0.616 ($\Delta = -0.163$). This is a substantial and consistent inflation.
    \item \textbf{Precision stays high while recall collapses.} RF temporal precision is 0.812 (even higher than random CV), but recall drops from 0.623 to 0.309---the model catches fewer than one in three post-2021 agents. This asymmetry means the classifier is conservative: when it flags an address as an agent, it is usually right, but it misses the majority of new-generation agents.
    \item \textbf{GBM is the most temporally robust model} by F1: temporal F1 of 0.604 vs.\ RF's 0.447 and LR's 0.537. GBM maintains better recall (0.499) than RF (0.309), sacrificing some precision.
    \item \textbf{Distribution shift drives the gap.} The training set is 70\% human (pre-2021 era, before agent platforms scaled), while the test set is 63\% agent (post-2021 era). The model learned a prior that ``most addresses are human,'' which fails in the agent-dominated test period. Additionally, gas-pricing features shift after EIP-1559 (August 2021), and newer agent platforms (Virtuals, AI Arena) use different behavioral patterns than older ones (Autonolas, Fetch.ai).
    \item \textbf{This finding strengthens the leakage narrative.} Random CV on provenance-verified labels (0.803) was already 18 points below the leaky AUC (0.98). Now we show that even the honest random-CV number is 13 points above what would be observed in a realistic deployment scenario where the model must generalize forward in time.
\end{enumerate}

\subsection{4.8 4-Dimensional Security Audit}

We run the 4-dimensional audit at three scales: the full 3,302-address leaky set, the n=549 provenance v3 subset (from the 563 set with 14 dropped for empty histories), and the trusted 50-address strict subset. Table 8b reports the n=549 provenance audit alongside the extremes.

\textbf{Table 8b. Security audit on n=549 provenance set (intermediate tier).}

\begin{table}[h]
\centering
\small
\begin{tabular}{lllll}
\toprule
Metric & Agent mean & Human mean & Ratio & p-value \\
\midrule
unlimited\_approval\_rate & 0.534 & 0.439 & \textbf{1.21x} & \textbf{0.0003} \\
dex\_interaction\_rate & 0.292 & 0.149 & \textbf{1.96x} & \textbf{<0.0001} \\
n\_unlimited\_approvals & 97.2 & 99.8 & 0.97x & 0.232 (NS) \\
revert\_rate & 0.041 & 0.042 & 0.97x & 0.701 (NS) \\
\bottomrule
\end{tabular}
\end{table}

On the intermediate n=549 set, agents have a significantly higher approval \textit{rate} (53.4\% vs 43.9\%) and DEX interaction rate (29.2\% vs 14.9\%), but similar raw approval \textit{counts} (97 vs 100). The extreme 10x ratio from the leaky set is gone; the honest finding is nuanced: agents approve more frequently but don't accumulate more approvals in absolute terms because humans hold approvals longer.

Table 8 below shows the leaky-vs-strict extremes.

\textbf{Table 8. 4-dimensional security audit: leaky vs honest.}

\begin{table}[h]
\centering
\small
\begin{tabular}{llllll}
\toprule
Metric (dimension) & Full 3,302 (C1-C4) agent / human & Agent/human ratio (leaky) & Trusted 50 agent / human & Agent/human ratio (honest) & Direction reversed? \\
\midrule
# approvals (D1) & 108.75 / 14.41 & \textbf{7.55x} & 110.09 / 125.11 & 0.88x & \textbf{Yes} \\
# unlimited approvals (D1) & 61.38 / 5.93 & \textbf{10.35x} & 49.14 / 66.14 & 0.74x & \textbf{Yes} \\
# active approvals (D1) & 102.02 / 13.49 & 7.56x & 104.68 / 121.14 & 0.86x & \textbf{Yes} \\
avg approval age (days, D1) & 1170.5 / 1134.9 & 1.03x & 795.6 / 1235.8 & 0.64x & \textbf{Yes} \\
max approval age (days, D1) & 1539.0 / 1295.5 & 1.19x & 1028.2 / 1811.9 & 0.57x & \textbf{Yes} \\
unlimited approval rate (D1) & 0.502 / 0.304 & 1.65x & 0.245 / 0.452 & 0.54x & \textbf{Yes} \\
DEX interaction rate (D3) & 0.228 / 0.133 & 1.71x & 0.071 / 0.041 & 1.73x & No \\
swap rate (D3) & 0.196 / 0.118 & 1.66x & 0.060 / 0.025 & 2.39x & No \\
revert rate (D4) & 0.049 / 0.042 & 1.16x & 0.015 / 0.024 & 0.64x & \textbf{Yes} \\
gas used efficiency (D4) & 0.712 / 0.695 & 1.02x & 0.644 / 0.751 & 0.86x & \textbf{Yes} \\
\bottomrule
\end{tabular}
\end{table}

\textbf{Key findings.}

1. \textbf{The agents-have-10x-more-approvals headline does not survive relabeling.} On the trusted set, humans hold \textbf{more} unlimited approvals per address than agents (mean 66.1 vs 49.1). This reversal is borderline significant (Mann-Whitney p=0.071 for unlimited approvals, p=0.026 for unlimited approval rate).
2. \textbf{Interpretation.} Famous ENS humans (vitalik.eth, Hayden Adams, nick.eth) are heavy long-term DeFi users with many historical approvals. The C1-C4 mining pipeline preferentially labeled token-holder addresses with heavy DeFi footprints as AGENT, which is exactly why the leaky ratio was 10x: we selected for agent-like humans and excluded human-like agents.
3. \textbf{DEX interaction rate stays directional} (1.7x on both sets) but the absolute magnitudes drop sharply (0.228 -> 0.071), indicating that the curated-MEV-bot subset is less DEX-active than the platform token holders.
4. \textbf{Revert rate reverses:} curated MEV bots have lower revert rates (1.5\%) than ENS humans (2.4\%), because MEV searchers simulate transactions before broadcasting.
5. \textbf{Gas used efficiency reverses:} curated humans are more efficient (0.75) than curated agents (0.64). This is the opposite of the intuitive claim that algorithms are more precise.

\textbf{The audit direction is not a minor wrinkle. It is a full 180-degree reversal on the headline "agents are riskier than humans on approvals".}

\subsection{4.9 Comparison with Published Bot-Detection Baselines}

To contextualize our honest numbers, Table~10 compares the provenance-verified AUC from this work with published bot-detection results across platforms.

\textbf{Table 10. Cross-platform bot-detection AUC comparison.}

\begin{table}[h]
\centering
\small
\begin{tabular}{lllll}
\toprule
Study & Platform & Task & Reported AUC & Label regime \\
\midrule
Cresci et al.\ (2017) & Twitter & Social bot detection & 0.75--0.85 & Crowd-sourced labels \\
Varol et al.\ (2017) & Reddit & Bot detection & $\sim$0.80 & Platform API flags \\
Chen et al.\ (2020) & Ethereum & Ponzi scheme detection & 0.92 & Heuristic + behavioral \\
Victor \& Weintraud (2021) & Ethereum & Wash-trade detection & 0.89 & Graph-rule labels \\
\midrule
\textbf{This work (leaky)} & Ethereum & Agent identification & 0.98 & C1-C4 heuristic \\
\textbf{This work (honest RF)} & Ethereum & Agent identification & \textbf{0.803} & Provenance-only \\
\textbf{This work (honest GAT)} & Ethereum & Agent identification & \textbf{0.832} & Provenance-only \\
\textbf{This work (temporal)} & Ethereum & Agent identification & 0.669 & Provenance + temporal split \\
\bottomrule
\end{tabular}
\end{table}

\textbf{Key findings.}

\begin{enumerate}
    \item \textbf{Our honest AUC (0.803--0.832) is competitive with state-of-the-art bot detection in the honest-label regime.} Twitter bot detection with crowd-sourced labels (Cresci et al.) achieves 0.75--0.85; Reddit bot detection (Varol et al.) reaches $\sim$0.80. Our honest numbers fall squarely in this range.
    \item \textbf{Higher published numbers likely reflect label leakage.} Chen et al.'s Ethereum Ponzi detector (AUC 0.92) uses behavioral features that overlap with their labeling heuristic, analogous to our C1-C4 leakage. Without an explicit leakage audit, such numbers cannot be compared directly to provenance-verified baselines.
    \item \textbf{The temporal holdout AUC (0.669) is the most realistic deployment number} and should be the reference for practitioners. The gap between 0.803 (random CV) and 0.669 (temporal holdout) quantifies the temporal optimism bias that affects all cross-validation-only evaluations.
    \item \textbf{Implication:} future bot-detection papers on blockchain data should report both random-CV and temporal-holdout numbers, along with an explicit leakage audit, before claiming to beat the 0.80 honest baseline.
\end{enumerate}


\section{5. Threats to Validity: The C1-C4 Label Leakage Discovery}

This section is the methodological centerpiece of the paper and documents the leakage discovery, the forensic evidence, the fix, and the implications.

\subsection{5.1 The Bug}

In our labeling script \texttt{paper1\_onchain\_agent\_id/features/verify\_c1c4.py}, the \texttt{\_check\_c3} function gates AGENT labels on three quantities:

- \texttt{tx\_interval\_cv} (coefficient of variation of inter-transaction intervals)
- \texttt{hour\_entropy} (Shannon entropy over UTC hour bucketing)
- \texttt{burst\_ratio} (fraction of consecutive transactions within 10 s)

These three quantities are algebraically the same as three of the top classifier features in the 23-feature vector:

- \texttt{tx\_interval\_std = tx\_interval\_mean x tx\_interval\_cv} (entangled with feature 2)
- \texttt{active\_hour\_entropy = hour\_entropy} (identical to feature 4)
- \texttt{burst\_frequency = burst\_ratio} (identical to feature 7)

Any address that was labeled AGENT because \texttt{hour\_entropy > 2.5} is then used to train a classifier whose top feature is \texttt{active\_hour\_entropy}. The model is \textbf{guaranteed} to re-derive the labeling rule from the features.

\subsection{5.2 The Evidence}

\textbf{Evidence 1: The honest vs leaky AUC gap.}

\begin{table}[h]
\centering
\small
\begin{tabular}{llll}
\toprule
Metric & Leaky (C1-C4, n=3316) & Full provenance (n=1,147) & Strict provenance (n=64) \\
\midrule
GBM AUC & 0.9803 & 0.8095 (10-fold) & 0.6276 (LOO) / 0.6940 (CV) \\
RF AUC & 0.9766 & \textbf{0.8028} (10-fold) & 0.7713 (LOO) / 0.8030 (CV) \\
GAT AUC & 0.9341 & \textbf{0.8318} (5-fold) & 0.8825 (5-fold) \\
Top feature univariate AUC & active\_hour\_entropy: 0.9347 & gas\_price\_trailing\_zeros\_mean: 0.1448 (importance) & tx\_interval\_mean: 0.7683 \\
\bottomrule
\end{tabular}
\end{table}

The AUC drops from 0.98 (leaky) to 0.80 (full provenance) to 0.77 (strict provenance) as we progressively tighten the labeling purity. The top-feature ranking also shifts: on the leaky set, the leaked temporal feature \texttt{active\_hour\_entropy} dominated; on n=1,147, gas features dominate; and on n=64, the temporal feature \texttt{tx\_interval\_mean} leads but with a much lower AUC (0.77 vs 0.93). This is the signature of label leakage.

\textbf{Evidence 2: Inter-label agreement is low and worsens at scale.}

We apply the GBM trained on the 64 provenance labels transductively to the 3,252 platform-token-holder rows (whose C1-C4 labels we no longer trust). The inter-label agreement between the honest prediction and the C1-C4 label is \textbf{37.4\%}, which is below random chance for balanced binary classification. The honest classifier does not endorse the majority of C1-C4 labels.

Furthermore, applying the C1-C4 heuristic directly to the 1,147 full provenance addresses yields a diagnostic agreement of only \textbf{48.8\%}, below chance. This was 55.8\% on the v3 set (n=563) and 92.9\% on the original 14-address v2 pilot, demonstrating that the C1-C4 heuristic's reliability degrades rapidly as the dataset diversifies beyond the narrow pilot categories. This is a sixth piece of evidence against heuristic labeling.

\textbf{Evidence 3: Statistical tests on n=64 reveal structural cracks.}

- RF significantly beats GBM (McNemar p=0.0225, DeLong p=0.0014, bootstrap CI95 [-0.24, -0.06]).
- GBM vs Majority: NOT significantly different (GBM LOO accuracy 0.578 vs Majority 0.516, McNemar p=0.58).
- Single best feature matches GBM.

These are the three symptoms you expect to see when a classifier's performance is being propped up by a leaking label: the strongest model (GBM) behaves like a majority baseline once the leak is sealed.

\textbf{Evidence 4: Cross-platform transfer reverses sign.}

Train on Autonolas C1-C4 labels and test on trusted provenance: AUC 0.24-0.34. Train on Fetch.ai C1-C4 labels and test on AI Arena C1-C4 labels: AUC 0.97-0.99. The gap is a direct measurement of the labeling-scheme-specific signal that the classifier picked up.

\textbf{Evidence 5: Security audit direction reverses.}

The "agents hold 10x more unlimited approvals" headline from the leaky set becomes "humans hold 1.3x more unlimited approvals" on the trusted set. The direction of a 4-dimensional audit is not supposed to depend on labeling heuristic. It does here because the heuristic preferentially admits heavy-DeFi users as AGENT.

\subsection{5.3 The Fix}

We decided on the following:

1. \textbf{Demote C1-C4 from a labeling oracle to a downstream diagnostic tool.} It can still be used to rank how likely an address is to be an agent given a trained classifier, but it cannot be the source of ground truth.
2. \textbf{Expand the provenance-labeled set from 64 to 1,147.} Through systematic v3 and v4 mining campaigns, we grew the honest corpus from 64 to 1,147 addresses across 18 provenance source categories, reducing the statistical uncertainty and enabling 10-fold CV instead of LOO.
3. \textbf{Retain n=64 as a strict sensitivity check.} The original 64 provenance-only addresses serve as a conservative reference point.
4. \textbf{Apply the RF trained on the 1,147 provenance labels transductively} to the 3,252 platform-token-holder rows to generate honest predictions, and publish these predictions as weak labels for future re-mining rather than as ground truth.
5. \textbf{Report leaky, full-provenance, and strict-provenance numbers} throughout the paper with clear labels.
6. \textbf{Publish the leakage case study} as Section 5 of this paper.

\subsection{5.4 A Reusable Leakage-Audit Toolkit}

The leakage we found is subtle: none of the classifier features are literally one of the gating quantities. The leak is through algebraic identities (\texttt{tx\_interval\_std $\approx$ tx\_interval\_cv $\times$ tx\_interval\_mean}) and name collisions (\texttt{hour\_entropy} gated the label, \texttt{active\_hour\_entropy} is a feature). A reviewer running a superficial check would not catch it; neither did we for several weeks.

We distill our experience into a three-step audit that any behavioral-classifier study on blockchain data can apply. Algorithm~\ref{alg:leakage\_audit} formalizes the procedure; if AUC drops monotonically across label-purity tiers \textbf{and} cross-scheme transfer collapses, leakage is confirmed.

\begin{algorithm}[t]
\SetAlgoLined
\KwIn{Feature set $\mathcal{F}$, labeling pipeline $\mathcal{L}$, classifier $h$}
\KwOut{Leakage verdict: \textsc{Confirmed} / \textsc{Not detected}}

\textbf{Step 1: Label--feature overlap check.}\;
Let $\mathcal{L}_q$ = all quantities used in $\mathcal{L}$ (thresholds, filters, derived scores)\;
\ForEach{$f \in \mathcal{F}$}{
    Compute Jaccard similarity $J(f, \ell)$ for each $\ell \in \mathcal{L}_q$, allowing algebraic transforms (log, ratio, CV $\leftrightarrow$ std/mean, renaming)\;
    \If{$J(f, \ell) > 0$ or $|\text{Pearson}(f, \ell)| > 0.9$}{
        Flag $(f, \ell)$ as entangled\;
    }
}
\If{no flags}{
    \Return \textsc{Not detected} (but proceed to Steps 2--3 as confirmation)\;
}

\BlankLine
\textbf{Step 2: AUC across label-purity tiers.}\;
Define tiers: $T_1$ = heuristic labels (leaky), $T_2$ = provenance labels, $T_3$ = strict provenance\;
Train $h$ on each tier; record AUC: $a_1, a_2, a_3$\;
\If{$a_1 > a_2 > a_3$ (monotonic decline with purity)}{
    Label-quality sensitivity confirmed\;
}

\BlankLine
\textbf{Step 3: Cross-label-scheme transfer.}\;
Train $h$ on $T_1$ labels, test on $T_2$ labels (and vice versa)\;
\If{cross-scheme AUC $< 0.5$ (worse than chance)}{
    The classifier learned the labeling scheme, not the concept\;
    \Return \textsc{Confirmed}\;
}
\Return \textsc{Not detected}\;
\caption{Three-Step Leakage Audit for Behavioral Classifiers}
\label{alg:leakage_audit}
\end{algorithm}

\textbf{Application to this study.} Step~1 identified three entangled pairs (\texttt{hour\_entropy}/\texttt{active\_hour\_entropy}, \texttt{burst\_ratio}/\texttt{burst\_frequency}, \texttt{tx\_interval\_cv}/\texttt{tx\_interval\_std}). Step~2 confirmed monotonic AUC decline: 0.98 (leaky) $\to$ 0.80 (full provenance) $\to$ 0.77 (strict provenance). Step~3 showed Autonolas$\to$trusted transfer AUC of 0.24--0.34, worse than chance. All three steps converge: leakage confirmed.

We release the implementing scripts (\texttt{verify\_c1c4.py}, \texttt{pipeline\_provenance.py}, \texttt{statistical\_tests.py}, \texttt{cross\_platform\_eval.py}) as a reusable toolkit that other researchers can adapt to their own labeling pipelines.

\subsection{5.5 Other Threats}

Beyond the C1-C4 issue, several additional threats apply.

- \textbf{\texttt{defi\_hf\_trader} label quality.} The largest agent category in the full provenance set is \texttt{defi\_hf\_trader} (199 addresses, 37\% of all agents), labeled as "agent" because the address makes >50 Uniswap router calls per active day. While this threshold strongly suggests automated execution, it is a behavioral criterion rather than a provenance-pure label. Some of these addresses could be humans using aggressive trading interfaces (e.g., custom front-ends with auto-routing). If a significant fraction are misclassified, the n=1,147 AUC would be artificially depressed (because the classifier is being penalized for correctly identifying humans that we mislabeled as agents) or inflated (if the behavioral criterion leaks into features the way C1-C4 did). We note that the >50-daily-calls threshold is at least one step removed from the 23 classifier features: none of the features directly encode daily Uniswap call count. However, \texttt{top\_contract\_concentration} and \texttt{method\_id\_diversity} are correlated with high-frequency Uniswap usage, so a residual leak is possible. This is the most important caveat on the n=1,147 results and should be investigated in future work by validating a random sample of \texttt{defi\_hf\_trader} addresses against deployment metadata or operator disclosures.
- \textbf{Sample size on the strict set.} n=64 is too small to support fine-grained claims. 5-fold CV has 13 test samples per fold; AUC confidence intervals are wide. The full provenance set (n=1,147) substantially mitigates this concern for the primary results.
- \textbf{Selection bias.} The strict set's 33 agents and 31 humans are public, named, famous addresses. The expanded set diversifies the sources but still over-represents high-frequency DeFi traders among agents and ENS-active users among humans.
- \textbf{Feature drift.} Agent behavior evolves as platforms update. Our features reflect a specific time window.
- \textbf{Adversarial evasion.} An adversary aware of the feature set can mimic human behavior by injecting jitter into gas pricing and transaction timing.
- \textbf{Chain coverage.} We only study Ethereum mainnet. Agents on Base, Arbitrum, Solana, and other chains are out of scope.
- \textbf{Contract-type confounding.} We excluded router contracts and exchange hot wallets from training, but edge cases remain.


\section{6. Discussion, Limitations, and Future Work}

\subsection{6.1 Practical Implications for Auditors and Regulators}

The temporal holdout results (Section~4.7) have direct consequences for how this classifier should be deployed in practice. The precision--recall asymmetry under temporal shift---precision stays high (0.81) while recall collapses (0.31)---means the following:

\begin{itemize}
    \item \textbf{High-confidence flagging works.} When the classifier labels an address as an agent, it is correct approximately 81\% of the time, even on post-2021 addresses that the model has never seen. This makes the classifier suitable as a \textbf{triage tool}: security auditors and monitoring services can use it to prioritize addresses for manual investigation, with a false-positive rate under 20\%.
    \item \textbf{Comprehensive detection does not work.} The classifier misses roughly 70\% of post-2021 agents (recall 0.31). New agent categories---LLM-driven strategy agents, AI16Z ELIZA deployments, Virtuals Protocol agents---exhibit behavioral patterns absent from the pre-2021 training distribution. The classifier should \textbf{not} be used as a definitive filter (e.g., to claim ``X\% of transactions are human''), because the false-negative rate is too high.
    \item \textbf{Periodic retraining is essential.} The 13-point AUC gap between random CV and temporal holdout quantifies the cost of stale features. As new agent platforms emerge and gas-pricing conventions evolve (e.g., EIP-4844 blob transactions, EIP-7702 delegation), the feature set must be updated and the model retrained. We recommend a 6-month retraining cadence with each cycle adding provenance-verified addresses from newly launched agent platforms.
    \item \textbf{Recommendation:} Deploy as a triage layer that surfaces high-confidence agent candidates for human review, not as an autonomous decision system. Pair with platform registry lookups and on-chain provenance checks for addresses that the classifier misses.
\end{itemize}

\subsection{6.2 Limitations}

1. \textbf{The full provenance corpus is moderately sized and heterogeneous.} At n=1,147 (533 agents, 614 humans), the full provenance set is large enough for 10-fold CV with reasonable fold sizes (~115 test samples per fold) and is approaching standard ML benchmarks. The 18 source categories vary in label quality: \texttt{defi\_hf\_trader} (199 addresses) uses a behavioral threshold (>50 daily Uniswap calls) that is one step removed from classifier features but not provenance-pure.
2. \textbf{The strict provenance corpus remains tiny.} The 64-address sensitivity check has wide uncertainty bands (5-fold CV std 0.08-0.12) and should not be over-interpreted.
3. \textbf{Source composition bias.} Agents are dominated by high-frequency DeFi traders (61\% of agent addresses), while humans are dominated by ENS-interaction-verified accounts (71\% of human addresses). This means the classifier is partly learning "high-frequency Uniswap trader vs casual ENS holder" rather than a general agent-vs-human distinction.
4. \textbf{The original C1-C4 mining exercise was not wasted, but it cannot be cited as ground truth.} The 3,252 platform-token-holder rows still serve as a test bed for feature engineering and as a large graph for GNN propagation, but their labels are now weak labels that must be cross-checked against provenance.
5. \textbf{GNN variance shrinks but is still meaningful.} GAT's 5-fold AUC std drops from 0.108 on n=64 to 0.021 on n=1,147, a substantial improvement. The n=1,147 GAT headline (0.832) is therefore statistically more reliable than the n=64 headline (0.883).
6. \textbf{The security audit cannot be generalized.} Even the honest direction is based on 22 agents and 28 humans. We cannot claim the direction reversal holds at scale; we can only claim that the 10x-ratio headline from the leaky set was an artifact.

\subsection{6.3 Clean Re-Mining Roadmap}

The v3 expanded mining partially fulfills the roadmap we originally proposed. Below we assess progress and remaining steps:

1. \textbf{Provenance-first labeling (substantially achieved).} The v4 mining added Compound V3 liquidators, Chainlink keepers, Gitcoin donors, ENS reverse setters, and PoolTogether depositors in addition to the v3 sources (Gelato registry, Keep3r bond registry, Flashbots builder list). However, the \texttt{defi\_hf\_trader} category uses a behavioral threshold (>50 daily Uniswap calls) that is not strictly provenance-based. Future work should validate this category against deployment metadata.
2. \textbf{Human controls at scale (achieved).} The v4 set includes 614 humans from 7 sources, up from 31 in the strict set. The human set now covers ENS-interaction-verified users (167), Gitcoin donors (146), ENS reverse setters (144), PoolTogether depositors (89), exchange depositors (28), expanded ENS-verified humans (25), and governance voters (5), substantially reducing the celebrity-human bias and achieving approximate class balance (0.87:1).
3. \textbf{Behavioral label audit (achieved).} We ran the 3-step leakage audit from Section 5.4 on all v3 candidate labels. The C1-C4 diagnostic agreement of 55.8\% confirms that the expanded labels are not simply recapitulating the heuristic.
4. \textbf{Held-out platform (not yet achieved).} All v3 sources contribute to training. A future version should reserve Gelato or Keep3r entirely for testing.
5. \textbf{Temporal holdout (achieved).} Section~4.7 reports a temporal split at the median first-seen block (April 2021). Random CV overestimates AUC by 13--16 points; recall collapses from 0.62 to 0.31 on post-2021 agents. Future work should explore rolling temporal splits at multiple cutoff dates.
6. \textbf{Adversarial stress tests (not yet achieved).} Inject behavioral noise into agent transactions to measure robustness against evasion.

\subsection{6.4 Longer-Term Directions}

- \textbf{EIP-7702 impact.} EIP-7702 lets an EOA temporarily delegate smart-contract code within a transaction, dissolving the \texttt{is\_contract()} test that has underpinned naive bot detection. Our behavioral framework does not depend on \texttt{is\_contract()}, so it should survive EIP-7702, but the set of viable features will change (e.g., authorization-list usage becomes a feature).
- \textbf{Cross-chain tracking.} Agents that run on multiple chains simultaneously (Base, Arbitrum, Solana) require cross-chain identity linking.
- \textbf{LLM-powered agents.} The companion Paper 3 (AI Sybil) work suggests that LLM-driven agents currently mimic rule-based agents in on-chain traces. We cannot yet separate LLM agents from rule-based agents from Etherscan data alone; new features and external signals are needed.
- \textbf{Agent-to-agent coordination detection.} The interaction-graph component of our GNN work hints that agent clusters are detectable, but the honest corpus is too small to make the claim.


\section{7. Conclusion}

We presented the first honest empirical study of AI agent identification on Ethereum mainnet using provenance-only ground truth. Our initial pipeline on 3,316 platform-mined addresses reported an apparently exciting AUC of 0.9803 (GBM). In the course of this work, we discovered that the labeling oracle used three quantities (\texttt{hour\_entropy}, \texttt{burst\_ratio}, \texttt{tx\_interval\_cv}) that were algebraically equivalent to three top classifier features, a direct form of target leakage. After demoting the oracle to a downstream diagnostic and assembling 1,147 provenance-verified addresses (533 agents from 11 source categories, 614 humans from 7 source categories), the honest baselines are:

\textbf{Primary results (n=1,147, full provenance):}

\begin{itemize}
    \item \textbf{Random Forest 10-fold AUC 0.803} (repeated 5x10-CV 0.803 $\pm$ 0.027),
    \item \textbf{Gradient Boosting 10-fold AUC 0.810},
    \item \textbf{GAT 5-fold AUC 0.832 $\pm$ 0.021} (the best honest model overall),
    \item \textbf{Temporal holdout AUC 0.669} (RF), revealing that random CV overestimates by 13 AUC points; precision stays high (0.81) but recall collapses to 0.31 on post-2021 agents.
\end{itemize}

\textbf{Sensitivity check (n=64, strict provenance):}

\begin{itemize}
    \item \textbf{Random Forest LOO-CV AUC 0.7713} (5x10-CV 0.8030),
    \item \textbf{GAT 5-fold AUC 0.8825} on the same subset (high variance: std 0.108),
    \item \textbf{GradientBoosting indistinguishable from a majority classifier} (McNemar p=0.58),
    \item \textbf{Autonolas $\to$ trusted cross-platform AUC 0.24--0.34}, worse than chance,
    \item \textbf{Security audit reverses direction}: humans hold more unlimited approvals than curated agents.
\end{itemize}

The four-tier evaluation (leaky 0.98 $\to$ full provenance 0.80 $\to$ strict provenance 0.77 $\to$ temporal holdout 0.67) demonstrates a progressive AUC decline as evaluation realism increases. The feature importance ranking also shifts: gas features dominate on n=1,147, displacing the temporal features that were inflated by the C1-C4 leak.

Beyond the honest baselines, we contribute a reusable three-step leakage-audit toolkit (Algorithm~\ref{alg:leakage\_audit}) that formalizes the checks we performed: (1)~label--feature overlap, (2)~AUC comparison across label-purity tiers, and (3)~cross-label-scheme transfer. If AUC drops monotonically across tiers and cross-scheme transfer collapses, leakage is confirmed. This toolkit is a standalone methodological contribution applicable to any behavioral classifier on blockchain data, and we release the implementing scripts alongside the data.

For practitioners, the precision--recall asymmetry under temporal shift provides clear deployment guidance: use the classifier as a high-confidence triage tool (81\% precision), not as a comprehensive filter (31\% recall). Periodic retraining with provenance-verified addresses from newly launched agent platforms is essential to maintain coverage as the agent population evolves.

On-chain AI agents are a real and growing phenomenon, and identifying them reliably is a real and growing need. Our central message is that this need will not be met by scaling up heuristic labeling pipelines. It will be met only by starting from non-behavioral provenance, running explicit leakage audits, reporting honest-vs-leaky gaps, and treating cross-label-scheme transfer as a non-negotiable diagnostic. The honest numbers are smaller than the leaky numbers, but they are the numbers the field can build on. All code, features, labels, and the leaky-vs-honest diagnostic artifacts are released.


\section{References}



\textit{This draft corresponds to the honest pipeline outputs of 2026-04-08 (commit eb8721e and follow-up). Headline numbers: RF LOO 0.7713 on n=64; GAT 5-fold 0.8825 on trusted subset; leaky-vs-honest gap of 0.21 AUC; security audit direction reversed. See \texttt{experiments/expanded/pipeline\_results\_provenance.json}, \texttt{statistical\_tests.json}, \texttt{cross\_platform\_eval.json}, \texttt{gnn\_results.json}, \texttt{combined\_pipeline\_results.json}, and \texttt{security\_audit\_expanded.json} for the raw JSON artifacts and \texttt{FINDINGS\_2026-04-08.md} for the session-level discovery log.}


\bibliographystyle{ACM-Reference-Format}
\bibliography{main}

\end{document}
